\section{隐私交易：交出 Context 与获得连接}

Elys 的产品难题，是让用户愿意交出 context。主持人指出，用户天然会有防御心理：不想让商业组织知道太多，不想让分身乱说话，也不确定交出 context 后有什么好处。Tristan 的回答是，产品必须在用户流失前让他感受到 context 飞轮的收益。

\begin{figure}[H]
\centering
\includegraphics[width=0.95\textwidth]{context-privacy-tradeoff.png}
\caption{Context 与隐私的交换：少交、适度交、过度交的收益与风险。}
\end{figure}

\begin{knowledgebox}{读图：隐私不是越多越好，context 也不是越多越好}
少交 context 风险低但产品无感；适度交 context 能带来更好匹配和连接；过度交 context 可能失控。好产品要让用户看到收益，并能控制授权范围。
\end{knowledgebox}

\subsection{本地化与信任}

访谈中提到，用户可能永远更信任本地保存的核心隐私。这意味着 AI 社交产品需要设计本地记忆、可撤回授权、透明日志、分级隐私和可解释分身行为。否则，context 越核心，用户越不敢交。

\subsection{真实连接 vs AI 热闹}

Tristan 明确说，AI 和 AI 之间互相发帖、互相评论意义不大。用户在意的是背后是否有真人。AI 可以先打前站、降低破冰门槛，但北极星应该是真人连接率和真人行为比例。

\begin{importantbox}{北极星指标：真人连接率}
AI 社交网络不是让 AI 自嗨，而是让真人更好连接。因此 Elys 的北极星不是 DAU 或 AI 消息数，而是真人连接率、真人行为比例和用户是否遇到真正有价值的人。
\end{importantbox}

\subsection{本章小结}

AI 社交必须处理 context 与隐私的交换。没有 context，分身无效；没有信任，用户不会交出 context；没有真人连接，AI 社交就失去意义。

